\section{Arithmetic reformulations and consequences}

The results of the preceding sections admit purely number-theoretic
formulations that do not refer to shift spaces, factor maps, or finite-state
presentations.  Throughout this section fix $m\ge 3$, $n\ge m$, and write
$L:=n+m-1$.

\subsection{Local Fibonacci congruence rigidity}

For a binary block $\omega=\omega_1\cdots\omega_L\in\{0,1\}^L$ and
$1\le t\le n$, define the $t$-th local Fibonacci weighted sum
\[
S_t^{(m)}(\omega):=\sum_{j=0}^{m-1}\omega_{t+j}F_{j+2},
\]
and its standard residue modulo $F_{m+2}$,
\[
R_t^{(m)}(\omega)\in\{0,1,\ldots,F_{m+2}-1\},
\qquad
R_t^{(m)}(\omega)\equiv S_t^{(m)}(\omega)\pmod{F_{m+2}}.
\]
Collect all local residues into a vector
\[
\mathcal R_{m,n}(\omega):=
\bigl(R_1^{(m)}(\omega),\ldots,R_n^{(m)}(\omega)\bigr)
\in\{0,1,\ldots,F_{m+2}-1\}^n.
\]

The key observation is that by Corollary~\ref{cor:visible-value-overflow},
\[
R_t^{(m)}(\omega)=V_m\bigl(\Fold_m(\omega_t\cdots\omega_{t+m-1})\bigr),
\]
where $V_m:X_m\to\{0,1,\ldots,F_{m+2}-1\}$ is the bijection from
Proposition~\ref{prop:visible-value-parametrization}.

\begin{theorem}[Local Fibonacci congruence rigidity]
\label{thm:congruence-rigidity}
The map
$\mathcal R_{m,n}:\{0,1\}^L\to\{0,1,\ldots,F_{m+2}-1\}^n$
is injective.  Equivalently, for any
$(r_1,\ldots,r_n)\in\{0,1,\ldots,F_{m+2}-1\}^n$,
the congruence system
\[
\sum_{j=0}^{m-1}\omega_{t+j}F_{j+2}\equiv r_t\pmod{F_{m+2}}
\qquad (1\le t\le n)
\]
has at most one solution $\omega\in\{0,1\}^L$.
The image $\mathcal R_{m,n}(\{0,1\}^L)$ has exactly $2^L=2^{n+m-1}$ elements.
\end{theorem}

\begin{proof}
Take any $\omega\in\{0,1\}^L$.  For each $t$, set
$a_t:=\Fold_m(\omega_t\cdots\omega_{t+m-1})\in X_m$.  By
Corollary~\ref{cor:visible-value-overflow},
$V_m(a_t)\equiv S_t^{(m)}(\omega)\pmod{F_{m+2}}$, and since $V_m(a_t)$
already lies in $\{0,1,\ldots,F_{m+2}-1\}$, one has
$V_m(a_t)=R_t^{(m)}(\omega)$.  Thus
\[
\mathcal R_{m,n}(\omega)=\bigl(V_m(a_1),\ldots,V_m(a_n)\bigr).
\]

Suppose $\mathcal R_{m,n}(\omega)=\mathcal R_{m,n}(\omega')$.  Then
$V_m(a_t)=V_m(a_t')$ for each $t$.  Since $V_m$ is a bijection
(Proposition~\ref{prop:visible-value-parametrization}), $a_t=a_t'$ for every
$1\le t\le n$.  Hence $\omega$ and $\omega'$ produce the same length-$n$
folded block under $B_{m,n}$.  By
Theorem~\ref{thm:block-bijection-threshold}, $B_{m,n}$ is a bijection, so
$\omega=\omega'$.  The image has $2^L$ elements because $\mathcal R_{m,n}$ is
an injection from a set of size $2^L$.
\end{proof}

\subsection{Fingerprint separation of equi-valued representations}

\begin{corollary}[Fingerprint separation]
\label{cor:fingerprint-separation}
For each integer $N$, let
\[
\Omega_L(N):=
\Bigl\{
\omega\in\{0,1\}^L:
\sum_{k=1}^L\omega_kF_{k+1}=N
\Bigr\}.
\]
Then $\mathcal R_{m,n}$ restricted to $\Omega_L(N)$ is still injective:
\[
\lvert\Omega_L(N)\rvert
=
\bigl\lvert\{\mathcal R_{m,n}(\omega):\omega\in\Omega_L(N)\}\bigr\rvert.
\]
\end{corollary}

\begin{proof}
Since $\mathcal R_{m,n}$ is injective on all of $\{0,1\}^L$
(Theorem~\ref{thm:congruence-rigidity}), it is injective on every subset.
\end{proof}

In other words, all length-$L$ binary Fibonacci representations of a single
integer $N$ have mutually distinct overlapping local Zeckendorf fingerprints.

\subsection{Exact bit recovery from residue windows}

Write
$\mathfrak R_{m,\ell}:=\mathcal R_{m,\ell}(\{0,1\}^{\ell+m-1})$
for the set of realizable local residue vectors of length $\ell$.

\begin{theorem}[Bit recovery from local congruence data]
\label{thm:bit-recovery-residues}
For each $0\le s\le m-1$, there exists a function
$\rho_{m,s}:\mathfrak R_{m,m}\to\{0,1\}$ such that for every
$\omega\in\{0,1\}^{n+m-1}$ and every $1\le t\le n-m+1$,
\[
\omega_{t+s}
=
\rho_{m,s}\bigl(
R_t^{(m)}(\omega),R_{t+1}^{(m)}(\omega),\ldots,R_{t+m-1}^{(m)}(\omega)
\bigr).
\]
In particular, taking $s=m-1$ gives a one-sided decoding rule
$g_m:=\rho_{m,m-1}$ with
\[
\omega_t=g_m\bigl(R_{t-m+1}^{(m)}(\omega),\ldots,R_t^{(m)}(\omega)\bigr).
\]
No rule using fewer than $m$ consecutive residues can recover all input bits.
\end{theorem}

\begin{proof}
For any admissible $m$-tuple $(r_1,\ldots,r_m)\in\mathfrak R_{m,m}$, set
$a_i:=V_m^{-1}(r_i)\in X_m$ and define
\[
\rho_{m,s}(r_1,\ldots,r_m):=\kappa_{m,s}(a_1\cdots a_m),
\]
where $\kappa_{m,s}$ is the local inverse rule from
Theorem~\ref{thm:Phi-m-conjugacy-threshold}.

Now take any $\omega\in\{0,1\}^{n+m-1}$ and fix $t$.  Set
$a_j:=\Fold_m(\omega_j\cdots\omega_{j+m-1})$ for $t\le j\le t+m-1$.  Then
$R_j^{(m)}(\omega)=V_m(a_j)$, so $a_j=V_m^{-1}(R_j^{(m)}(\omega))$.  By
Theorem~\ref{thm:Phi-m-conjugacy-threshold}, the folded block
$a_t\cdots a_{t+m-1}$ has unique raw lift
$\omega_t\cdots\omega_{t+2m-2}$, and $\kappa_{m,s}$ reads its $(s+1)$-st
bit.  Therefore
\[
\rho_{m,s}\bigl(R_t^{(m)}(\omega),\ldots,R_{t+m-1}^{(m)}(\omega)\bigr)
=
\kappa_{m,s}(a_t\cdots a_{t+m-1})
=
\omega_{t+s}.
\]

For the optimality claim, suppose a rule using only $k+1<m$ consecutive
residues existed.  Since $V_m$ is a symbol-by-symbol bijection, this would
yield a sliding-block left inverse $\Psi:Y_m\to\{0,1\}^{\ZZ}$ with memory
$k<m-1$, contradicting
Corollary~\ref{cor:Phi-m-memory-threshold-exact}.
\end{proof}

\subsection{Distribution of random local residue vectors}

\begin{theorem}[Exact distribution of random local residues]
\label{thm:random-residue-distribution}
Let $X_1,\ldots,X_L$ be independent with
$\mathbb P(X_i=1)=p$ and $\mathbb P(X_i=0)=1-p$,
where $0<p<1$.  Define the random local residue vector
$\mathbf R_{m,n}:=\mathcal R_{m,n}(X_1\cdots X_L)$.  Then for every
$\mathbf r=(r_1,\ldots,r_n)\in\{0,1,\ldots,F_{m+2}-1\}^n$,
\[
\mathbb P(\mathbf R_{m,n}=\mathbf r)=
\begin{cases}
p^{|\omega(\mathbf r)|_1}(1-p)^{L-|\omega(\mathbf r)|_1},
& \mathbf r\in\mathfrak R_{m,n},\\[4pt]
0,& \mathbf r\notin\mathfrak R_{m,n},
\end{cases}
\]
where $\omega(\mathbf r)$ is the unique preimage guaranteed by
Theorem~\textup{\ref{thm:congruence-rigidity}}.  When $p=\tfrac12$,
$\mathbf R_{m,n}$ is uniformly distributed on $\mathfrak R_{m,n}$:
\[
\mathbb P(\mathbf R_{m,n}=\mathbf r)=2^{-L}
\qquad (\mathbf r\in\mathfrak R_{m,n}).
\]
\end{theorem}

\begin{proof}
If $\mathbf r\notin\mathfrak R_{m,n}$, then no
$\omega\in\{0,1\}^L$ satisfies $\mathcal R_{m,n}(\omega)=\mathbf r$, so the
probability is $0$.

Suppose $\mathbf r\in\mathfrak R_{m,n}$.  Set
$a_t:=V_m^{-1}(r_t)\in X_m$ for $1\le t\le n$.  The event
$\{\mathbf R_{m,n}=\mathbf r\}$ is equivalent to the folded block being
$a_1\cdots a_n$.  By Theorem~\ref{thm:block-bijection-threshold}, this
folded block has a unique raw lift $\omega(\mathbf r)\in\{0,1\}^L$.  Hence
\[
\{\mathbf R_{m,n}=\mathbf r\}
=
\{(X_1,\ldots,X_L)=\omega(\mathbf r)\}.
\]
Since $X_1,\ldots,X_L$ are independent Bernoulli$(p)$,
\[
\mathbb P(\mathbf R_{m,n}=\mathbf r)
=
p^{|\omega(\mathbf r)|_1}(1-p)^{L-|\omega(\mathbf r)|_1}.
\]
When $p=\tfrac12$, every $\omega\in\{0,1\}^L$ has probability $2^{-L}$.
\end{proof}

\begin{corollary}[Entropy of random local residues]
\label{cor:entropy-random-residues}
Under the assumptions of Theorem~\textup{\ref{thm:random-residue-distribution}},
for every $n\ge m$,
\[
H(\mathbf R_{m,n})=(n+m-1)h_2(p),
\]
where $h_2(p):=-p\log_2 p-(1-p)\log_2(1-p)$.  In particular, when
$p=\tfrac12$, $H(\mathbf R_{m,n})=n+m-1$.
\end{corollary}

\begin{proof}
By Theorem~\ref{thm:random-residue-distribution}, the distribution of
$\mathbf R_{m,n}$ is the image of a length-$L$ Bernoulli block under the
bijection $\mathcal R_{m,n}$.  Entropy is invariant under bijective relabeling,
so $H(\mathbf R_{m,n})=H(X_1,\ldots,X_L)=Lh_2(p)=(n+m-1)h_2(p)$.
\end{proof}

\subsection{Markov order of the residue process}

\begin{theorem}[Exact Markov order of the local residue process]
\label{thm:residue-markov-order}
Let $(X_t)_{t\in\ZZ}$ be a two-sided Bernoulli$(p)$ process with $0<p<1$.
Define
\[
R_t:=
\biggl(\sum_{j=0}^{m-1}X_{t+j}F_{j+2}\biggr)\bmod F_{m+2},
\qquad
R_t\in\{0,1,\ldots,F_{m+2}-1\}.
\]
Then $(R_t)_{t\in\ZZ}$ is an exact $m$-step Markov process: given any
realizable block $(r_{t-m+1},\ldots,r_t)$, the conditional law of $R_{t+1}$
depends only on these $m$ values, and no $(m-1)$-step Markov law can describe
$(R_t)$.

More precisely, let $a_i:=V_m^{-1}(r_{t-m+i})\in X_m$,
$\beta:=a_1\cdots a_m\in\mathcal L_m(Y_m)$, and let
$v(\beta)\in\{0,1\}^{m-1}$ be the terminal suffix of the unique raw lift of
$\beta$.  Then for each $b\in\{0,1\}$,
\[
\mathbb P\!\left(
R_{t+1}=V_m(\Fold_m(v(\beta)b))
\;\middle\vert\;
R_{t-m+1}=r_{t-m+1},\ldots,R_t=r_t
\right)
=
p^b(1-p)^{1-b}.
\]
\end{theorem}

\begin{proof}
Set $Y_t:=\Fold_m(X_t\cdots X_{t+m-1})\in X_m$.  By
Corollary~\ref{cor:visible-value-overflow} and
Proposition~\ref{prop:visible-value-parametrization}, $R_t=V_m(Y_t)$ for all $t$.
Since $V_m$ is a bijection, $(R_t)$ and $(Y_t)$ are related by a
symbol-by-symbol relabeling.

Theorem~\ref{thm:exact-markov-order-bernoulli} states that $(Y_t)$ is an exact
$m$-step Markov process under $\mu_{m,p}$, with the displayed transition
probabilities.  Translating through $V_m$ yields the same conclusion for
$(R_t)$: the event $\{Y_j=a_{j-(t-m)}\}$ is equivalent to $\{R_j=r_j\}$
because $R_j=V_m(Y_j)$ and $V_m$ is invertible.  Likewise
$Y_{t+1}=\Fold_m(v(\beta)b)$ if and only if
$R_{t+1}=V_m(\Fold_m(v(\beta)b))$.

If $(R_t)$ were $(m-1)$-step Markov, then $(Y_t)=V_m^{-1}(R_t)$ would also be
$(m-1)$-step Markov, contradicting
Theorem~\ref{thm:exact-markov-order-bernoulli}.
\end{proof}

